\section{Capacity toward the heavy targets}
\label{sec:minimum-heavy-capacity}

Let $D$ be strongly connected, all-deficient, and minimal under
deletion of arbitrary nonempty sets of arcs.  Assume
$n=2\delta+3$ and $\delta\ge6$.
Let $K$ be its positive-surplus core, let $\mathcal H$ be the set of
heavy vertices ($\sigma\ge2$), and put $H=|\mathcal H|$.
The established general result gives $H\ge2$.
Use the certified basins $B_1,\ldots,B_H$ of
Lemma~\ref{lem:general-basin-certificates}, and the exceptional set
\[
 B=\{x\in K:e(x)=0,\ \sigma(x)=1,\ p(x)=0\}.
\]
No source is far from a vertex of $B$.  Every other core vertex
has a heavy-root certificate.  A source pointing to a root points
to its entire basin; a source reaching a root within two steps
reaches its entire basin within two steps.
Also, every minimum-degree source has at least two far vertices
in $K$, all certified: its far set has size three or four and
contains at most one zero-surplus target.
Every heavy root has a minimum-degree predecessor.

\begin{lemma}[Minimum-degree capacity toward the heavy set]
\label{lem:minimum-heavy-capacity}
For every minimum-degree vertex $v$,
\[
 {|N_D^+(v)\cap\mathcal H|\le H-2.}
\]
\end{lemma}
\begin{proof}
If $v$ is itself a heavy root and points to every other heavy
root, it points to its own children and to all the other basins.
It would have no far certified core target, a contradiction.
Thus consider a minimum-degree vertex $v$ that is not heavy.
Suppose it points to at least $H-1$ heavy roots.
It cannot point to all $H$, since that would directly reach every
certified core target.  Let $r$ be the unique heavy root not
dominated by $v$.

Every far core target of $v$ is in the basin of $r$.
The certificate therefore makes $r$ far from $v$.
Choose any minimum-degree predecessor $w\to r$.
Vertex $w$ points to the basin of $r$ and has a far certified
core target elsewhere, hence some root $s\ne r$ is far from $w$.
As $v\to s$, neither $v\to w$ nor $w\to v$ is possible: the
former would let $v$ reach $r$, and the latter would let $w$
reach $s$ within two steps.

If $v$ is whole, this nonadjacency is impossible.
If $\mu(v)=1$, the same reasoning for every minimum-degree
predecessor of $r$ confines that entire predecessor set to the
unique nonneighbor of $v$.  It is nonempty, so $c(r)=1$;
its member is nonwhole, so $c_0(r)=0$.
The refined predecessor bound gives
\[
 9c(r)+6c_0(r)\ge\delta+1+3\Phi(r)\ge\delta+4,
\]
contradicting $\delta\ge6$.

These cases cover every minimum-degree vertex outside $K$ and
every inactive minimum-degree unit target.
In the remaining case $v$ is an active minimum-degree unit
target and has a heavy-root certificate.  Its certified root
cannot be one of the roots it dominates, so it must be $r$.
The containment $N^-(r)\subseteq N^-(v)$ makes $w\to v$,
contradicting the preceding nonadjacency (equivalently,
$w\to v\to s$ reaches a far root).
All possible types are excluded.
\end{proof}

\begin{theorem}[Capacity for every vertex]
\label{thm:all-heavy-capacity}
Every vertex $v$ of $D$ satisfies
\begin{equation}\label{eq:all-heavy-capacity}
 {|N_D^+(v)\cap\mathcal H|\le H-2.}
\end{equation}
In particular $H\ge3$.  If $\delta\ge10$, then $H\ge4$ and
\[
 {z\ge h+10,}
\]
where $z$ counts whole minimum-degree vertices and $h$ counts
gap-two vertices.
\end{theorem}
\begin{proof}
Minimum-degree vertices were handled above.  If a heavy vertex
$v$ points to all the other heavy roots, a minimum predecessor
$w\to v$ reaches its own root directly and every other root
within two steps through $v$.  It therefore reaches every
certified core target within two steps, contradicting the
minimum-source far demand.

The remaining vertices are nonminimum unit targets, and all
have a heavy-root certificate, say with root $r$.
The arc $r\to v$ prevents $v$ from pointing back to $r$.
If it pointed to every other heavy root, any minimum predecessor
$w\to r$ would point to $v$ by
$N^-(r)\subseteq N^-(v)$, and would again reach every root and
certified core target within two steps.  This contradiction
proves \eqref{eq:all-heavy-capacity}.

If $H=2$, the minimum-degree capacity is zero, contrary to the
existence of a minimum-degree predecessor for every heavy root.
Thus $H\ge3$.
If $H=3$, delete the set $S$ of all three heavy roots.
Each surviving minimum-degree vertex loses at most one outgoing
arc by \eqref{eq:all-heavy-capacity}.
Every surviving nonminimum vertex is unit and has excess one,
so its outdegree after deletion is also at least $\delta-1$.
All excess-two vertices are heavy and have been deleted.
Hence $\delta^+(D-S)\ge\delta-1$.
The heavy-triple obstruction in
Corollary~\ref{cor:heavy-triple} contradicts $\delta\ge10$.
Consequently $H\ge4$ in that range.
The general core bound gives $k\ge6$, while the fixed-target
budget gives
$z-h=\sum_x\sigma(x)\ge k+H$.
This yields $z\ge h+10$.
\end{proof}

\begin{corollary}[A pair covering when there are four heavy roots]
\label{cor:four-heavy-minimum-pairs}
Suppose $\delta\ge10$ and $H=4$.  Form a simple graph $J$ on
$\mathcal H$, joining two heavy roots if some minimum-degree
vertex outside that pair points to both.  Every three-element
subset of $\mathcal H$ contains an edge of $J$; equivalently,
the independence number of $J$ is at most two.
More explicitly, with
$a_v=|N_D^+(v)\cap\mathcal H|$, one has
\[
 4\le\sum_{e(v)=0,\ a_v=2}\bigl(2-[v\in\mathcal H]\bigr).
\]
\end{corollary}
\begin{proof}
The heavy-triple obstruction requires an outside minimum-degree
vertex pointing to at least two members of each triple, or an
outside excess-one vertex pointing to all three.
The latter alternative is ruled out by
\eqref{eq:all-heavy-capacity}, which limits every heavy outdegree
to two.  Thus every triple contains a pair dominated by a
minimum-degree source outside that triple.
Such a source has $a_v=2$ and covers two triples if it is not
heavy, and one triple if it is itself heavy; summation gives
the explicit inequality.
\end{proof}
